\documentclass[border=2mm]{standalone}

\newif\ifset \settrue

\ifset
\usepackage{pgfmath}
\else
\usepackage{tikz}
\fi




\begin{document}

% >>> texdoc tikz p998
% \pgfmathparse{<函数,方法>}\let\mydef=\pgfmathresult
% \usepackage{pgfmath} \usepackage{tikz}



% 设置一个宏,注意这是静态
\pgfmathsetmacro{\dx}{rand*0.1}
\dx,\dx,\dx,









% 解算宏内容
\pgfmathparse{1 ? 4 : 2}\let\mymax\pgfmathresult
\mymax





% 用\def的 可以展开命令
% 额外命令,可以先设定好随机种子
\pgfmathsetseed{10}
\def\dy{\pgfmathparse{rand*0.1}\pgfmathresult}
\dy,\dy,\dy

% 用\def的 使用 noexpand 设定为不展开命令
\def\dy{\noexpand\pgfmathparse{max(3,4)}\pgfmathresult}
\dy 




\pgfmathparse{rand}\pgfmathresult						% [-1,-1] 随机数
\pgfmathparse{random()}\pgfmathresult					% [0,1] 随机数
\pgfmathparse{random(100)}\pgfmathresult				% 小于100的正随机数
\pgfmathparse{random(232,762)}\pgfmathresult			% 区间中的随机数

\pgfmathparse{div(75,9)} \pgfmathresult					% 去掉余数,返回正整数
\pgfmathparse{sqrt(8765.432)} \pgfmathresult			% 根号
\pgfmathparse{pow(2,7)} \pgfmathresult					% 
\pgfmathparse{abs(-5)} \pgfmathresult					% 取到绝对值
\pgfmathparse{mod(20,6)} \pgfmathresult					% 余数2
\pgfmathparse{sign(-5)} \pgfmathresult					% 判断正负数
\pgfmathparse{round(32.5/17)} \pgfmathresult			% 4舍5入 取整数
\pgfmathparse{floor(-398/12)} \pgfmathresult			% 向下取整数
\pgfmathparse{ceil(32.5/17)} \pgfmathresult				% 向上取整数
\pgfmathparse{int(32.5/17)} \pgfmathresult				% 取整数
\pgfmathparse{frac(32.5/17)} \pgfmathresult 			% 商值取小数
\pgfmathparse{real(4)} \pgfmathresult					% 转为小数点
\pgfmathparse{veclen(12,5)} \pgfmathresult		% 两数平方根
\pgfmathparse{gcd(42,56)} \pgfmathresult		% 两数相减





\pgfmathparse{min(3,4,-2,250,-8,100)} \pgfmathresult	% min函数
\pgfmathparse{max(3,4,-2,250,-8,100)} \pgfmathresult	% max函数

\pgfmathparse{not(true)} \pgfmathresult					% not函数
\pgfmathparse{ifthenelse(5==4,"yes","no")} \pgfmathresult
\pgfmathparse{true ? "yes" : "no"} \pgfmathresult
\pgfmathparse{false ? "yes" : "no"} \pgfmathresult		% 二元判断语句


\pgfmathparse{iseven(2)} \pgfmathresult			% 是否为偶数
\pgfmathparse{isodd(2)} \pgfmathresult			% 是否为奇数
\pgfmathparse{isprime(1)} \pgfmathresult		% 是否素数



\pgfmathparse{greater(20,25)} \pgfmathresult	% 大于
\pgfmathparse{less(20,25)} \pgfmathresult
\pgfmathparse{notequal(20,25)} \pgfmathresult	% 不等于
\pgfmathparse{notgreater(20,25)} \pgfmathresult	% 大于等于
\pgfmathparse{notless(20,25)} \pgfmathresult	% 小于等于
\pgfmathparse{and(5>4,6>7)} \pgfmathresult		% 并条件
\pgfmathparse{or(5>4,6>7)} \pgfmathresult		% 或条件
    


%文本
\pgfmathparse{depth("Some Lovely Text")} \pgfmathresult		% 获得文本框深度
\pgfmathparse{height("Some Lovely Text")} \pgfmathresult	% 获得文本框高度
\pgfmathparse{width("Some Lovely Text")} \pgfmathresult		% 获得文本框宽度






\end{document}